\section{Le dimensionnement à l'épreuve des faits}
\label{sec:sizing}

Le moteur ne choisit pas seulement \emph{quand} être investi, mais \emph{avec
quelle taille}. Cette taille est censée être pilotée par la volatilité~: plus le
marché s'agite, plus la position se réduit, de façon à viser une volatilité
constante de 55\,\% annualisés. Ce réglage a été calibré sur l'or. La figure
suivante montre ce qu'il devient sur une paire de change majeure.

\begin{figure}[H]
\centering
\includegraphics[width=\linewidth]{usdjpy_trades_sizing}
\caption{En haut, le levier retenu à l'entrée de chaque trade et son plafond.
En bas, la volatilité réalisée d'USD/JPY sur 21~séances, annualisée, et le
seuil sous lequel le plafond de levier devient contraignant.}
\end{figure}

\begin{insightbox}
\textbf{Sur le change, c'est le plafond qui décide, pas la cible.} La
volatilité réalisée d'USD/JPY évolue entre 2,9 et 18,1\,\% sur la période,
médiane \metric{8,2\,\%} --- très loin de la cible de 55\,\%. À cette
volatilité médiane, le mécanisme demanderait un levier de 6,7$\times$ qu'il
n'obtient pas~: le
plafond de 6,6$\times$ mord dès que la volatilité passe sous
\metric{8,3\,\%}, c'est-à-dire une séance sur deux et pour \textbf{18~entrées
sur 35}.
\end{insightbox}

Ce constat ressemble à celui de la sleeve or, mais il faut le lire plus
finement. Le plafond mord souvent~; il ne mord pas fort. À la volatilité
médiane du yen, $6{,}6 \times 8{,}2 = 54\,\%$~: la position délivre presque
exactement la volatilité visée. Ce n'est que dans les épisodes de calme extrême
--- fin 2021, la volatilité tombe à 2,9\,\% --- que la sleeve tourne loin sous
son objectif, autour de 19\,\% de volatilité effective. Autrement dit, le
couple \emph{(cible 55\,\%, plafond 6,6$\times$)}, calibré pour l'or, se trouve
presque juste sur USD/JPY par coïncidence arithmétique.

Cette coïncidence est un point de fragilité, pas un argument. Elle signifie que
le comportement du dimensionnement sur ce candidat dépend de deux paramètres qui
n'ont jamais été réglés pour lui~: relever la cible n'aurait presque aucun
effet tant que le plafond ne bouge pas, et relever le plafond augmenterait
directement le risque dans les phases calmes, c'est-à-dire précisément là où le
signal se trompe le plus souvent.

\subsection{Ce que le pas de lot impose --- et ce qu'il n'impose pas}

Le simulateur ne peut pas exprimer un levier arbitraire~: il achète des lots
par pas de 0,01. Sur la sleeve or, ce poste était un candidat sérieux pour
expliquer les écarts entre moteurs, parce que les volumes y sont minuscules ---
entre 0,09 et 0,26~lot, soit un pas qui représente jusqu'à 11\,\% de la
position.

Ici, ce n'est pas le cas. Les volumes vont de 7,85 à 24,60~lots~: le pas de
0,01 représente entre \metric{0,04 et 0,13\,\%} de la position. L'écart entre
le levier voulu et le levier obtenu est donc négligeable, et ce poste peut être
écarté sans regret --- ce qui est une bonne nouvelle méthodologique~: sur ce
candidat, l'écart éventuel entre recherche et exécution ne pourra pas être
imputé à l'arrondi des volumes.

\begin{warningbox}
Le dimensionnement observé ici est celui d'un \textbf{run de recherche}~:
sleeve seule, allocation 100\,\%, facteur de risque 1,0, dépôt 100\,000\,\$. En
portefeuille, la sleeve or tourne à 10\,\% d'allocation avec un facteur de
risque de 4,5, et le cycle a documenté qu'un facteur de risque non épinglé dans
un run isolé ruine le compte. Les leviers et les montants de ce chapitre ne se
transposent donc pas~: une intégration éventuelle de ce candidat imposerait de
refaire le dimensionnement de bout en bout.
\end{warningbox}
